% =====================================================================
%  Chapter 1 — Introduction  (~1,500 words)
%  DRAFT generated 2026-07-18. Re-read in your own voice before submission;
%  every number is either a \macro from numbers.tex or marked with a comment.
% =====================================================================
\chapter{Introduction}
\label{ch:intro}

\section{Motivation}
\label{sec:intro-motivation}

Retrieval-augmented generation (RAG) has become the default engineering answer
to a well-known failure of large language models: they hallucinate. Ground the
model in retrieved documents, the argument runs, and its answers inherit the
reliability of the source material. In consumer applications this is a
reasonable default. In medicine it is not, for three reasons that together
motivate this dissertation.

The first is \emph{safety}. A hallucinated drug interaction is not merely a
wrong answer; in a deployed clinical tool it is a potential clinical incident.
Closed-book generation --- answering purely from the model's parametric
knowledge --- is therefore never a neutral choice in this domain: it trades
latency for an unbounded and invisible hallucination risk, with no provenance
trail a clinician could audit. Any system intended even for clinical
\emph{education} must take a position on when parametric knowledge alone is
trustworthy.

The second is \emph{cost}. On the hardware this project targets, retrieval is
not free. A retrieving query on the CPU-only reference configuration takes on
the order of \latPoneMcq{} of wall-clock time, dominated by the long-context
generation over the retrieved passages; a closed-book answer to the same
question takes a fraction of that. Multiplied over every query in a working
day, always-retrieve is a substantial and often unnecessary expenditure of
time and energy.

The third is \emph{deployment reality}. The intended setting for this work is
offline, CPU-only hardware: hospital training wards, rural clinics, hospitals
in low- and middle-income countries, and air-gapped environments where data
governance forbids sending clinical text to a third-party API. There is no
cloud model to call and no GPU to accelerate retrieval-augmented inference.
Every design decision in this project --- a 4-bit quantised 3B-parameter
model, exact CPU-side vector search, energy profiling --- follows from taking
this constraint seriously rather than treating it as an afterthought.

These three pressures converge on a single question that is usually skipped:
\emph{retrieval is neither free nor always helpful --- so when should a
clinical question-answering system actually retrieve?}

\section{Problem Statement}
\label{sec:intro-problem}

The two degenerate retrieval policies bracket the problem, and both fail.

\emph{Always-retrieve} maximises cost, and --- as this dissertation
demonstrates empirically --- it can also \emph{reduce} accuracy. On
multiple-choice medical questions, injecting retrieved context broke 46
previously-correct answers and fixed only 13 (Chapter~\ref{ch:eval}):
retrieval actively distracts a model that already knew the answer. The
assumption that retrieved context is at worst neutral does not survive contact
with the data.

\emph{Closed-book} minimises cost but offers no provenance and no bound on
hallucination. Its failures are silent: the model produces a fluent answer
whether or not the underlying knowledge exists, and a reader has no signal to
distinguish the two cases.

The gap between these extremes is a \emph{per-query} decision: retrieve
exactly when the model does not already know. Making that decision in this
deployment setting carries two hard constraints. It must be
\emph{training-free}, because no labelled ``should-retrieve'' dataset exists
for medicine and fine-tuning the generator (as reflection-token approaches
require) is infeasible offline on CPU. And it must require \emph{no second
model}, because the compute budget is a single quantised 3B model on a
laptop. The remaining option --- the one this dissertation develops and
evaluates --- is to gate retrieval on signals the model already emits during
generation: the entropy of its token distributions, the margin between its top
candidates, and the stability of its answer under reframing.

\section{Research Questions}
\label{sec:intro-rqs}

The evaluation in Chapter~\ref{ch:eval} is organised around four research
questions, each of which the collected data can actually answer:

\begin{description}
  \item[RQ1] How do always-retrieve, hybrid, closed-book and gated retrieval
    policies compare on accuracy, retrieval cost and latency for medical
    question answering on CPU-only hardware?
  \item[RQ2] Do training-free gating signals (token entropy, logit margin,
    answer-stability probing) and their published operating thresholds
    transfer from the general-domain literature to a small quantised medical
    model?
  \item[RQ3] Does an ensemble of three signals outperform its individual
    members, and if so, why?
  \item[RQ4] How far is the best achievable policy from clinically acceptable
    accuracy at each level of question risk --- the \emph{safety envelope}?
\end{description}

% NOTE: the original plan's hardware-tier RQ is deliberately absent — tiers
% were implemented but not physically benchmarked; see §6.2.

\section{Contributions}
\label{sec:intro-contributions}

This dissertation makes six contributions, each anchored to specific evidence:

\begin{enumerate}
  \item \textbf{A training-free retrieval-gating ensemble.} Three signals ---
    mean token entropy, mean top-two logit margin, and a two-draft
    hallucination probe --- combined by majority vote, requiring no
    fine-tuning, no labelled gate data, and no auxiliary model
    (Chapter~\ref{ch:method}).
  \item \textbf{Calibration non-transfer.} Literature-default thresholds
    ($\tau_H = 2.5$ nats, $\tau_M = 0.3$) are shown to be \emph{physically
    unreachable} on a 4-bit quantised 3B model, whose observed entropy spans
    0.27--1.59 nats and margin 0.45--0.89. Gate thresholds must be
    recalibrated per model; an offline replay procedure does so in seconds
    (Section~\ref{sec:eval-calibration}).
  \item \textbf{Retrieval as distraction.} On multiple-choice questions over a
    426K-chunk medical corpus, retrieval breaks 46 previously-correct answers
    and fixes 13 (net $-16.5$pp). Retrieved chunks are equally on-topic in
    both groups but $2.4\times$ less answer-bearing where retrieval harms:
    the failure mode is distraction by relevant-but-not-decisive context, not
    retriever error (Section~\ref{sec:eval-distraction}).
  \item \textbf{The protective skip.} On exactly those retrieval-poisoned
    questions, the gate's \textsc{skip} decision scored 19/19 (100\%); its
    \textsc{retrieve} decision, 3/27 (11\%). The gate is protective precisely
    where retrieval is harmful (Section~\ref{sec:eval-protective}).
  \item \textbf{The safety envelope.} A risk-stratified analysis quantifying
    the \emph{distance} between this system and clinical acceptability: the
    best policy on the benchmark falls \gapToLowBar{}pp short of the
    \barLow{} low-risk bar. The contribution is an honest measurement of that
    distance, not a claim of clinical readiness
    (Section~\ref{sec:eval-envelope}).
  \item \textbf{A reproducible evaluation harness.} Every reported number is
    generated from raw run logs into \LaTeX{} macros, guarded by a regression
    test that fails the build on a hand-edited literal; 180+ unit tests run
    without the model via a mock backend; runs are resumable and
    deterministic (Chapter~\ref{ch:impl}).
\end{enumerate}

\section{Dissertation Structure}
\label{sec:intro-structure}

Chapter~\ref{ch:background} situates the work among adaptive-retrieval
methods, medical RAG benchmarks and the uncertainty-signal literature, and
states the four-part gap this project occupies. Chapter~\ref{ch:method}
presents the system design: the policy space, the retrieval stack, and the
three-gate ensemble with its hardware-aware degradation.
Chapter~\ref{ch:impl} documents the implementation, with particular attention
to the reproducibility machinery. Chapter~\ref{ch:eval} reports the
evaluation: the main policy comparison, the calibration study, the
distraction and protective-skip analyses, the gate ablation, token-level
entropy attribution, and the safety envelope. Chapter~\ref{ch:discussion}
synthesises the findings and confronts their limitations.
Chapter~\ref{ch:lsep} addresses the legal, social, ethical and professional
dimensions of deploying (or declining to deploy) such a system.
Chapter~\ref{ch:conclusion} concludes and orders the future work by value.
